%bismi ALLAH ARRAHMAN ARRAHIM


















































\documentclass[11pt,a4paper]{article}

% ============================================================
% PACKAGES
% ============================================================

\usepackage[utf8]{inputenc}
\usepackage[T1]{fontenc}
\usepackage[french]{babel}
\usepackage{amsmath,amssymb,amsthm,mathtools}
\usepackage{bm}
\usepackage{geometry}
\usepackage{enumitem}
\usepackage{booktabs}
\usepackage{array}
\usepackage{hyperref}
\usepackage{microtype}

\geometry{margin=2.6cm}

% ============================================================
% THEOREM ENVIRONMENTS
% ============================================================

\newtheorem{definition}{Définition}[section]
\newtheorem{hypothesis}[definition]{Hypothèse}
\newtheorem{assumption}[definition]{Postulat de modélisation}
\newtheorem{proposition}[definition]{Proposition}
\newtheorem{theorem}[definition]{Théorème}
\newtheorem{lemma}[definition]{Lemme}
\newtheorem{corollary}[definition]{Corollaire}
\newtheorem{remark}[definition]{Remarque}

% ============================================================
% COMMANDS
% ============================================================

\newcommand{\X}{\mathcal X}
\newcommand{\D}{\mathcal D}
\newcommand{\R}{\mathbb R}
\newcommand{\N}{\mathbb N}
\newcommand{\Z}{\mathbb Z}
\newcommand{\HH}{\mathcal H}
\newcommand{\W}{\mathcal W}
\newcommand{\V}{\mathcal V}
\newcommand{\E}{\mathcal E}
\newcommand{\F}{\mathcal F}
\newcommand{\I}{\mathcal I}
\newcommand{\T}{\mathcal T}
\newcommand{\Rel}{\mathcal R}
\newcommand{\supp}{\operatorname{supp}}
\newcommand{\rank}{\operatorname{rank}}
\newcommand{\argmax}{\operatorname*{arg\,max}}
\newcommand{\argmin}{\operatorname*{arg\,min}}

% ============================================================
% DOCUMENT
% ============================================================

\begin{document}

\begin{center}
    {\LARGE \textbf{Théorie des empreintes et architecture
    informationnelle des modèles de langage}}\\[0.5cm]
    {\large Formalisation du modèle observationnel, mémoriel,
    relationnel et émergent de la représentation distribuée}
\end{center}

\vspace{0.5cm}

\begin{abstract}
Ce document formalise une hypothèse théorique selon laquelle la compréhension
de la mémoire, de la connaissance, de la généralisation et de certains aspects
de l'intelligence d'un modèle de langage doit être recherchée non pas
prioritairement dans des unités localisées représentant individuellement des
concepts, mais dans l'organisation distribuée des structures informationnelles
produites par l'apprentissage.

Le cadre proposé est construit à partir de la théorie des empreintes
informationnelles. L'observation constitue l'opération fondamentale par
laquelle un système accède à une partie d'un domaine d'information. Une
empreinte est alors une représentation produite relativement à un dispositif
d'observation, à une résolution et à un contexte. Plusieurs observations
peuvent produire des empreintes partiellement communes. Leur organisation
forme un maillage relationnel dans lequel une même composante interne peut
participer à plusieurs structures sémantiques distinctes.

Dans cette perspective, un neurone poly-sémantique n'est pas nécessairement
un neurone qui ``contient'' plusieurs concepts incompatibles. Il peut être une
composante commune à plusieurs empreintes distribuées. La généralisation est
alors interprétée comme l'exploitation de structures communes à plusieurs
empreintes, tandis que la spécialisation correspond à l'apparition de
distinctions supplémentaires lorsque la résolution représentationnelle
augmente.

Le présent texte distingue explicitement les définitions mathématiques, les
conséquences formelles, les hypothèses de modélisation et les conjectures
empiriques. Il ne prétend donc pas démontrer que cette interprétation est
déjà établie expérimentalement ; il fournit le formalisme nécessaire pour la
rendre testable.
\end{abstract}

\tableofcontents

% ============================================================
\section{Objet et statut de la théorie}
% ============================================================

\subsection{Problème général}

Les modèles de langage de grande taille (LLM) présentent des comportements
qui semblent difficilement réductibles à une interprétation locale de leurs
unités internes. En particulier, des analyses de neurones peuvent montrer
qu'une même unité est activée dans des contextes sémantiques apparemment très
différents.

Un exemple conceptuel est constitué par une unité activée à la fois dans des
contextes relatifs à un \emph{chat} et à une \emph{voiture}. Une interprétation
naïve consisterait à attribuer au neurone une représentation composite du type
``chat + voiture''. Une autre hypothèse est cependant possible : le neurone
pourrait être une composante commune de plusieurs structures distribuées dont
les significations globales sont différentes.

La présente théorie étudie cette seconde possibilité.

\subsection{Principe méthodologique}

Le principe méthodologique fondamental est le suivant :

\begin{quote}
Une unité interne ne doit pas être identifiée avec la signification globale
d'une représentation sans avoir déterminé la structure relationnelle dans
laquelle cette unité intervient.
\end{quote}

Ainsi, on distingue systématiquement :

\[
\text{unité interne}
\neq
\text{empreinte}
\neq
\text{concept}
\neq
\text{connaissance}.
\]

Une unité peut participer à plusieurs empreintes sans que ces empreintes
soient identiques.

\subsection{Statut épistémique}

Le formalisme contient quatre niveaux qu'il convient de ne pas confondre :

\begin{enumerate}[label=\arabic*.]
    \item les définitions mathématiques ;
    \item les propriétés déductibles de ces définitions ;
    \item les hypothèses de modélisation appliquées aux LLM ;
    \item les hypothèses empiriques qui doivent être vérifiées
    expérimentalement.
\end{enumerate}

Une conséquence mathématique du formalisme ne constitue donc pas
automatiquement une affirmation empirique sur les LLM.

% ============================================================
\section{Cadre observationnel général}
% ============================================================

\subsection{Domaine des objets}

Soit $\X$ un domaine d'objets, d'états, d'événements ou de configurations
susceptibles d'être observés.

Une observation n'est jamais considérée ici comme un accès absolu à l'objet.
Elle est définie relativement à un dispositif observationnel.

On introduit :

\[
\Phi :
\mathcal P \times \X \times \F
\longrightarrow
\R,
\]

où :

\begin{itemize}
    \item $\mathcal P$ est l'espace des contextes ou configurations
    observationnelles ;
    \item $\F$ est une famille de sondes ou fonctions d'observation ;
    \item $\Phi(P,x,f)$ est la réponse obtenue lorsqu'on observe $x$
    dans le contexte $P$ au moyen de la sonde $f$.
\end{itemize}

\subsection{Empreinte}

\begin{definition}[Empreinte observationnelle]
Pour $P\in\mathcal P$, $\lambda\subseteq\F$ et $x\in\X$, on définit
l'empreinte de $x$ relativement à $(P,\lambda)$ par

\[
I_{\Phi,\lambda}^{P}(x)
=
\left(
f\mapsto \Phi(P,x,f)
\right)_{f\in\lambda}.
\]

L'empreinte est donc une fonction sur l'ensemble des observations disponibles.
\end{definition}

Cette définition implique immédiatement que deux objets peuvent être
distincts tout en possédant la même empreinte relativement à une famille
d'observations donnée.

\begin{definition}[Indiscernabilité observationnelle]
Pour $x,y\in\X$, on pose

\[
x\sim_{\Phi,P,\lambda}y
\quad\Longleftrightarrow\quad
I_{\Phi,\lambda}^{P}(x)
=
I_{\Phi,\lambda}^{P}(y).
\]

\end{definition}

La relation $\sim_{\Phi,P,\lambda}$ est une relation d'équivalence.

Elle définit donc une partition de $\X$ en classes d'indiscernabilité.

\subsection{Résolution}

Une famille de sondes peut être organisée par un paramètre de résolution
$\lambda$ appartenant à un ensemble ordonné $(\Lambda,\leq)$.

On suppose que

\[
\lambda_1\leq\lambda_2
\]

signifie que $\lambda_2$ permet au moins autant de distinctions
observationnelles que $\lambda_1$.

On obtient alors une famille :

\[
\left(
\sim_{\lambda}
\right)_{\lambda\in\Lambda}.
\]

Lorsque la résolution augmente, les classes d'équivalence peuvent se
subdiviser :

\[
[x]_{\lambda_2}
\subseteq
[x]_{\lambda_1}
\qquad
\text{si }
\lambda_1\leq\lambda_2.
\]

Ainsi :

\[
\boxed{
\text{augmentation de résolution}
\Longrightarrow
\text{augmentation potentielle du nombre de distinctions}.
}
\]

% ============================================================
\section{Mémoire, connaissance et intelligence}
% ============================================================

\subsection{Principe général}

Dans le cadre présent, il est préférable de ne pas identifier la mémoire à un
conteneur d'objets et la connaissance à une collection d'objets stockés.

On introduit plutôt trois notions structurelles.

\begin{definition}[Mémoire]
Une mémoire est une structure capable de conserver, de manière persistante,
des traces ou des transformations d'empreintes produites par des observations
antérieures.
\end{definition}

Symboliquement :

\[
\boxed{
\text{Mémoire}
=
\text{persistance structurée d'empreintes ou de transformations
d'empreintes}.
}
\]

\begin{definition}[Connaissance]
Une connaissance est une organisation relationnelle suffisamment stable
d'un ensemble d'empreintes pour permettre leur identification, comparaison,
transformation ou réutilisation dans un contexte ultérieur.
\end{definition}

Ainsi :

\[
\boxed{
\text{Connaissance}
=
\text{organisation relationnelle des empreintes}.
}
\]

\begin{definition}[Intelligence]
Dans le présent cadre, l'intelligence désigne la capacité dynamique d'un
système à exploiter, transformer, combiner ou sélectionner des structures
issues de ses empreintes afin de produire des distinctions, prédictions,
décisions ou transformations adaptées à un problème.
\end{definition}

Ainsi :

\[
\boxed{
\text{Intelligence}
=
\text{exploitation dynamique des structures d'empreintes}.
}
\]

Ces trois notions ne sont donc pas indépendantes :

\[
\boxed{
\text{observation}
\rightarrow
\text{empreinte}
\rightarrow
\text{mémoire}
\rightarrow
\text{organisation}
\rightarrow
\text{connaissance}
\rightarrow
\text{exploitation dynamique}
\rightarrow
\text{intelligence}.
}
\]

% ============================================================
\section{L'Univers comme domaine informationnel}
% ============================================================

\subsection{Hypothèse philosophique et modèle théorique}

Une intuition fondamentale du modèle est que l'Univers peut être considéré,
au niveau informationnel, comme un domaine extrêmement vaste de phénomènes
susceptibles de produire des traces observables.

Cette proposition doit être distinguée de l'affirmation physique selon
laquelle l'Univers serait littéralement une mémoire.

On adopte donc la formulation plus faible :

\begin{assumption}[Domaine informationnel]
L'Univers peut être modélisé, pour les besoins de cette théorie, comme un
domaine $\mathcal U$ contenant des états et des phénomènes dont certaines
propriétés peuvent produire des empreintes accessibles à des observateurs.
\end{assumption}

Ainsi :

\[
\mathcal U
\supseteq
\X_{\mathrm{obs}},
\]

où $\X_{\mathrm{obs}}$ désigne l'ensemble des phénomènes effectivement
accessibles à un système donné.

\subsection{Limitation de l'observateur}

Un observateur individuel ne peut généralement pas accéder à l'ensemble
de $\mathcal U$.

Il dispose d'un ensemble limité de sondes :

\[
\lambda_O\subsetneq\F.
\]

Son état informationnel est donc :

\[
I_{\lambda_O}(x),
\]

et non une représentation absolue de $x$.

La connaissance acquise par un observateur est par conséquent une projection
partielle du domaine informationnel accessible.

\subsection{Accumulation collective}

Soient $O_1,\ldots,O_m$ plusieurs observateurs possédant des familles de
sondes :

\[
\lambda_{O_1},\ldots,\lambda_{O_m}.
\]

Une coopération peut produire :

\[
\lambda_{\mathrm{collectif}}
=
\bigcup_{i=1}^{m}\lambda_{O_i}.
\]

Sous l'hypothèse que ces observations sont complémentaires, on peut avoir :

\[
\lambda_{\mathrm{collectif}}
\supsetneq
\lambda_{O_i}.
\]

La résolution collective peut donc dépasser celle d'un observateur isolé.

On obtient alors conceptuellement :

\[
\boxed{
\text{coopération}
\rightarrow
\text{élargissement des observations}
\rightarrow
\text{augmentation de résolution}
\rightarrow
\text{nouvelles distinctions}
\rightarrow
\text{nouvelles structures de connaissance}.
}
\]

Il s'agit ici d'une conséquence du modèle observationnel, et non d'une loi
empirique universelle sur les sociétés humaines.

% ============================================================
\section{Passage aux modèles de langage}
% ============================================================

\subsection{Dataset}

Soit

\[
\D=\{d_1,d_2,\ldots,d_N\}
\]

un dataset d'apprentissage.

Chaque donnée $d_i$ est une séquence de tokens :

\[
d_i=(t_{i,1},t_{i,2},\ldots,t_{i,L_i}).
\]

On note $\mathcal V$ le vocabulaire du modèle :

\[
t_{i,j}\in\mathcal V.
\]

Le dataset constitue l'environnement informationnel auquel le modèle est
exposé pendant l'apprentissage.

\subsection{Représentation vectorielle}

À chaque token $t\in\mathcal V$ est associé une représentation initiale :

\[
e(t)\in\R^{d}.
\]

Cependant, la représentation effectivement utilisée par les couches internes
dépend du contexte.

On écrit donc :

\[
h_{i,j}^{(0)}
=
E(t_{i,j}),
\]

puis, pour les couches $\ell=1,\ldots,L$ :

\[
h_{i,j}^{(\ell)}
=
F_{\ell}
\left(
h_{i,1}^{(\ell-1)},
\ldots,
h_{i,L_i}^{(\ell-1)}
\right).
\]

La représentation finale est :

\[
h_{i,j}^{(L)}.
\]

Ainsi, un token n'est pas nécessairement associé à un vecteur invariant
indépendant du contexte.

% ============================================================
\section{Empreinte interne d'un LLM}
% ============================================================

\subsection{Définition}

Soit un modèle paramétré par $\theta$.

Pour une donnée $d$, on définit son état interne comme :

\[
H_\theta(d)
=
\left\{
h_{j}^{(\ell)}
\right\}_{j,\ell}.
\]

On appelle cette structure une \emph{empreinte interne} de la donnée
relativement au modèle $\theta$ :

\[
\boxed{
I_\theta(d)=H_\theta(d).
}
\]

Cette définition ne prétend pas que l'état interne constitue l'unique
empreinte pertinente. On peut également définir des projections :

\[
I_{\theta,\lambda}(d)
=
P_\lambda(H_\theta(d)),
\]

où $P_\lambda$ sélectionne une famille d'observables internes correspondant
à une résolution $\lambda$.

\subsection{Empreinte d'un token contextualisé}

Pour le token $t_j$ dans le contexte $d$, on peut définir :

\[
I_\theta(t_j\mid d)
=
\left(
h_j^{(0)},h_j^{(1)},\ldots,h_j^{(L)}
\right).
\]

Ainsi, le même token lexical peut produire différentes empreintes :

\[
I_\theta(t\mid d_1)
\neq
I_\theta(t\mid d_2).
\]

La signification n'est donc pas nécessairement localisée dans le token
lui-même.

% ============================================================
\section{Dépendance relationnelle et ``récursion''}
% ============================================================

\subsection{Problème terminologique}

L'intuition initiale emploie le terme ``récursion'' pour désigner le fait
qu'une donnée provoque des dépendances entre plusieurs dimensions, plusieurs
tokens et plusieurs niveaux de représentation.

Le terme ``récursion'' est trop large pour constituer directement une
définition mathématique.

On introduit donc une notion plus précise :

\begin{definition}[Dépendance relationnelle itérée]
Une dépendance relationnelle itérée est une chaîne de transformations dans
laquelle l'état représentatif à un niveau donné dépend d'états antérieurs,
du contexte et des interactions entre plusieurs composantes.
\end{definition}

Formellement :

\[
H^{(\ell+1)}
=
F_\ell
\left(
H^{(\ell)},C^{(\ell)}
\right),
\]

où $C^{(\ell)}$ représente l'information contextuelle accessible au niveau
$\ell$.

La répétition :

\[
H^{(0)}
\rightarrow
H^{(1)}
\rightarrow
\cdots
\rightarrow
H^{(L)}
\]

constitue le processus relationnel itéré.

\subsection{Structure relationnelle d'une donnée}

À une donnée $d_i$ est associée une structure :

\[
R_i
=
R_\theta(d_i).
\]

Cette structure peut contenir :

\begin{itemize}
    \item les tokens ;
    \item leurs positions ;
    \item leurs représentations ;
    \item les dépendances entre tokens ;
    \item les interactions entre dimensions ;
    \item les activations des unités internes ;
    \item les transformations produites par les couches.
\end{itemize}

Ainsi :

\[
d_i
\longrightarrow
R_i.
\]

Le dataset produit donc une famille :

\[
\mathcal R_\theta(\D)
=
\{R_1,\ldots,R_N\}.
\]

% ============================================================
\section{Maillage informationnel}
% ============================================================

\subsection{Définition}

Les structures $R_i$ ne sont généralement pas supposées indépendantes.

Deux données distinctes peuvent partager une composante interne :

\[
R_i\cap R_j\neq\varnothing.
\]

On définit alors le maillage informationnel du modèle par :

\[
\boxed{
\mathfrak M_\theta(\D)
=
\bigcup_{i=1}^{N}R_\theta(d_i),
}
\]

muni des relations d'intersection, de dépendance et de transformation
entre ses composantes.

Le maillage n'est donc pas une simple somme des données.

Il résulte du recouvrement des structures produites par leur traitement
commun.

\subsection{Superposition}

Soient deux données $d_a$ et $d_b$ telles que :

\[
R_a\neq R_b.
\]

Il peut néanmoins exister :

\[
R_a\cap R_b\neq\varnothing.
\]

On appelle ce phénomène \emph{superposition structurelle}.

Il ne signifie pas :

\[
d_a=d_b
\]

ni :

\[
\operatorname{sem}(d_a)=\operatorname{sem}(d_b).
\]

Il signifie seulement que certaines composantes de leurs structures
internes sont communes.

% ============================================================
\section{Neurone comme composante structurelle}
% ============================================================

\subsection{Unité interne}

Soit $\mathcal N=\{n_1,\ldots,n_M\}$ l'ensemble des unités internes étudiées.

Pour chaque unité $n_k$, on définit une fonction d'activation :

\[
a_k:
\D
\rightarrow
\R.
\]

Une unité est donc une fonction qui participe aux états internes produits
par le traitement des données.

\subsection{Support d'activation}

Pour un seuil $\tau\geq0$, on définit :

\[
\supp_\tau(n_k)
=
\{d\in\D:
|a_k(d)|\geq\tau\}.
\]

Une unité peut donc intervenir dans plusieurs familles de données :

\[
d_a,d_b\in\supp_\tau(n_k).
\]

Cela ne permet cependant pas de conclure que l'unité représente
individuellement la signification commune aux données.

\subsection{Hypothèse de compositionalité distribuée}

\begin{hypothesis}[Compositionalité distribuée]
Une empreinte interne sémantiquement significative est généralement une
structure distribuée de plusieurs unités, et non une unité isolée.
\end{hypothesis}

On écrit :

\[
I_\theta(d)
\approx
\mathcal C_\theta(d)
\subseteq
\mathcal N,
\]

où $\mathcal C_\theta(d)$ désigne une structure de composantes internes
participant conjointement à l'empreinte de $d$.

Cette hypothèse est compatible avec l'existence d'unités poly-sémantiques.

% ============================================================
\section{Poly-sémanticité}
% ============================================================

\subsection{Définition}

Soit $d_a$ un contexte relatif à un chat et $d_b$ un contexte relatif à une
voiture.

Supposons :

\[
a_k(d_a)\neq0
\]

et :

\[
a_k(d_b)\neq0.
\]

On pourrait conclure naïvement :

\[
n_k
=
\text{représentation de ``chat + voiture''}.
\]

La présente théorie refuse cette inférence.

\subsection{Interprétation structurelle}

Supposons plutôt :

\[
\mathcal C_\theta(d_a)
\subseteq\mathcal N
\]

et

\[
\mathcal C_\theta(d_b)
\subseteq\mathcal N.
\]

Il est possible que :

\[
\mathcal C_\theta(d_a)
\cap
\mathcal C_\theta(d_b)
\neq
\varnothing.
\]

Si :

\[
n_k
\in
\mathcal C_\theta(d_a)
\cap
\mathcal C_\theta(d_b),
\]

alors $n_k$ participe aux deux empreintes.

On obtient :

\[
\boxed{
n_k\in
\mathcal C_\theta(d_a)
\cap
\mathcal C_\theta(d_b)
\quad\not\Rightarrow\quad
n_k=\operatorname{concept}(d_a,d_b).
}
\]

Le neurone est alors un élément commun d'une structure distribuée.

% ============================================================
\section{Exemple : chat et voiture}
% ============================================================

Considérons :

\[
d_c=\text{``Le chat mange.''}
\]

et :

\[
d_v=\text{``La voiture roule.''}
\]

On suppose :

\[
\mathcal C_\theta(d_c)
=
\{n_1,n_4,n_8,n_{17},n_{31}\}
\]

et :

\[
\mathcal C_\theta(d_v)
=
\{n_2,n_8,n_{17},n_{24},n_{31}\}.
\]

Alors :

\[
\mathcal C_\theta(d_c)
\cap
\mathcal C_\theta(d_v)
=
\{n_8,n_{17},n_{31}\}.
\]

Il existe donc une superposition.

Mais :

\[
d_c\neq d_v.
\]

La signification globale des deux empreintes reste distincte.

La théorie propose alors de rechercher les propriétés relationnelles
communes responsables de l'intersection.

Par exemple, de manière hypothétique :

\[
\text{entité}
\rightarrow
\text{action}
\rightarrow
\text{état dynamique}.
\]

Cette interprétation n'est toutefois qu'une hypothèse tant qu'elle n'est pas
validée par une analyse expérimentale des représentations.

% ============================================================
\section{Généralisation}
% ============================================================

\subsection{Principe}

Soient plusieurs données :

\[
d_1,\ldots,d_m
\]

dont les empreintes sont différentes :

\[
I_\theta(d_i)\neq I_\theta(d_j)
\qquad(i\neq j).
\]

Supposons néanmoins qu'elles partagent une structure :

\[
S
\subseteq
\bigcap_{i=1}^{m}
\mathcal C_\theta(d_i).
\]

Alors $S$ constitue une structure commune aux différentes empreintes.

\begin{definition}[Structure généralisable]
Une structure $S$ est dite généralisable relativement à une famille
$\{d_1,\ldots,d_m\}$ si elle est identifiable dans leurs empreintes
respectives et demeure exploitable pour une nouvelle donnée présentant une
structure compatible.
\end{definition}

La généralisation peut alors être représentée schématiquement par :

\[
\boxed{
\text{plusieurs empreintes}
\rightarrow
\text{structure commune}
\rightarrow
\text{réutilisation}
\rightarrow
\text{généralisation}.
}
\]

La généralisation ne nécessite donc pas que les données soient identiques.

Elle nécessite l'existence d'une structure commune suffisamment stable.

% ============================================================
\section{Résolution et abstraction}
% ============================================================

\subsection{Résolution interne}

On introduit une famille de projections :

\[
P_\lambda:
\mathcal H
\rightarrow
\mathcal H_\lambda,
\]

où $\lambda$ représente un niveau de résolution.

L'empreinte observée à cette résolution est :

\[
I_{\theta,\lambda}(d)
=
P_\lambda(I_\theta(d)).
\]

Deux données peuvent être indiscernables à faible résolution :

\[
I_{\theta,\lambda_1}(d_a)
=
I_{\theta,\lambda_1}(d_b),
\]

mais distinguables à résolution supérieure :

\[
I_{\theta,\lambda_2}(d_a)
\neq
I_{\theta,\lambda_2}(d_b),
\qquad
\lambda_1<\lambda_2.
\]

On retrouve ainsi le principe fondamental de la théorie des empreintes :

\[
\boxed{
\text{résolution faible}
\rightarrow
\text{moins de distinctions},
}
\]

et :

\[
\boxed{
\text{résolution élevée}
\rightarrow
\text{davantage de distinctions}.
}
\]

\subsection{Abstraction}

Une abstraction peut être interprétée comme une structure qui reste stable
lorsque certaines distinctions sont supprimées.

Soit :

\[
P_{\lambda_1}
\]

une projection de résolution inférieure à :

\[
P_{\lambda_2}.
\]

Une structure $S$ est abstraite relativement à
$(\lambda_1,\lambda_2)$ si :

\[
P_{\lambda_1}(S)
\]

conserve une propriété relationnelle présente dans :

\[
P_{\lambda_2}(S).
\]

Ainsi, une abstraction peut être conçue comme une structure stable à
plusieurs niveaux de résolution.

% ============================================================
\section{Lien avec l'indiscernabilité de la théorie des empreintes}
% ============================================================

On définit :

\[
d_1
\sim_{\theta,\lambda}
d_2
\quad\Longleftrightarrow\quad
I_{\theta,\lambda}(d_1)
=
I_{\theta,\lambda}(d_2).
\]

Cette relation définit une partition du domaine des données.

Lorsque la résolution augmente, on peut obtenir :

\[
[d]_{\lambda_2}
\subseteq
[d]_{\lambda_1}.
\]

Le système passe donc d'une représentation grossière :

\[
\X/{\sim_{\lambda_1}}
\]

à une représentation plus fine :

\[
\X/{\sim_{\lambda_2}}.
\]

La hiérarchie des représentations peut ainsi être décrite par :

\[
\X/{\sim_{\lambda_1}}
\longrightarrow
\X/{\sim_{\lambda_2}}
\longrightarrow
\cdots.
\]

Cette structure fournit une interprétation formelle de l'idée selon laquelle
une augmentation de résolution permet de raffiner progressivement la
connaissance.

% ============================================================
\section{Seuil de distinction dans un LLM}
% ============================================================

Pour deux données $x,y$, on peut définir un seuil de distinction interne :

\[
\Theta_\theta(x,y)
=
\sup
\left\{
\iota(\lambda):
x\sim_{\theta,\lambda}y
\right\},
\]

sous réserve des conditions nécessaires sur l'échelle $\iota$.

Cette quantité mesure le niveau maximal auquel les deux données restent
indiscernables relativement à la famille d'observations choisie.

On obtient alors :

\[
\Theta_\theta(x,y)
\]

comme analogue interne du seuil observationnel de la théorie des empreintes.

Sous les hypothèses appropriées sur la filtration des observations, on
retrouve une structure ultramétrique potentielle :

\[
\Theta_\theta(x,z)
\geq
\min
\{
\Theta_\theta(x,y),
\Theta_\theta(y,z)
\}.
\]

Cette propriété ne doit toutefois être revendiquée que si les hypothèses
permettant la construction de la filtration sont effectivement satisfaites
dans le système étudié.

% ============================================================
\section{Du dataset au maillage de connaissances}
% ============================================================

Le processus d'apprentissage est schématisé par :

\[
\D
\xrightarrow{\text{observation}}
\mathfrak M_\theta(\D)
\xrightarrow{\text{optimisation}}
\theta^\ast.
\]

Plus explicitement :

\[
d_i
\longrightarrow
I_{\theta}(d_i)
\longrightarrow
R_\theta(d_i)
\]

pour chaque :

\[
i=1,\ldots,N.
\]

L'ensemble des structures produit :

\[
\mathfrak M_\theta(\D)
=
\bigcup_{i=1}^{N}R_\theta(d_i).
\]

L'apprentissage modifie progressivement $\theta$, donc les structures
produites par les données.

On peut ainsi représenter :

\[
\theta_0
\rightarrow
\theta_1
\rightarrow
\cdots
\rightarrow
\theta_T
=
\theta^\ast.
\]

À chaque étape :

\[
\mathfrak M_{\theta_t}(\D)
\]

peut évoluer.

% ============================================================
\section{Apprentissage comme construction d'un espace relationnel}
% ============================================================

Soit une fonction de perte :

\[
\mathcal L(\theta;\D).
\]

L'apprentissage cherche :

\[
\theta^\ast
\in
\argmin_\theta
\mathcal L(\theta;\D).
\]

La présente théorie ajoute une interprétation structurelle :

\[
\boxed{
\text{optimisation}
\rightarrow
\text{modification des représentations}
\rightarrow
\text{modification du maillage d'empreintes}.
}
\]

Ainsi, le dataset n'est pas simplement mémorisé comme une liste.

Il impose des contraintes statistiques et relationnelles sur l'espace des
paramètres.

Le modèle cherche une organisation de $\theta$ permettant de réduire la
perte sur l'ensemble des observations d'apprentissage.

% ============================================================
\section{Mémoire distribuée}
% ============================================================

\subsection{Principe}

Si une information donnée est distribuée sur plusieurs paramètres et
plusieurs unités, il n'est pas nécessaire qu'un emplacement unique
corresponde à cette information.

On peut représenter une information $K$ par :

\[
K
\rightsquigarrow
\mathcal S_K
\subseteq
\mathcal N\times\mathcal L,
\]

où $\mathcal S_K$ est son support structurel.

Deux connaissances $K_1,K_2$ peuvent donc satisfaire :

\[
\mathcal S_{K_1}
\cap
\mathcal S_{K_2}
\neq
\varnothing.
\]

Cette propriété fournit un modèle formel de la mémoire distribuée et de la
superposition.

\subsection{Conséquence}

Le fait qu'une composante participe à deux connaissances ne signifie pas que
ces connaissances sont identiques :

\[
\boxed{
\mathcal S_{K_1}\cap\mathcal S_{K_2}\neq\varnothing
\not\Rightarrow
K_1=K_2.
}
\]

Il s'agit précisément du mécanisme permettant théoriquement de comprendre
la poly-sémanticité sans attribuer plusieurs concepts contradictoires à une
même unité isolée.

% ============================================================
\section{Connaissance comme graphe d'empreintes}
% ============================================================

On peut représenter le système par un graphe :

\[
G_\theta=(V,E),
\]

où les sommets peuvent représenter :

\begin{itemize}
    \item des tokens ;
    \item des représentations ;
    \item des unités internes ;
    \item des empreintes ;
    \item des structures relationnelles.
\end{itemize}

Les arêtes représentent des relations telles que :

\begin{itemize}
    \item dépendance ;
    \item coactivation ;
    \item transformation ;
    \item attention ;
    \item projection ;
    \item interaction ;
    \item partage structurel.
\end{itemize}

Une connaissance n'est alors plus nécessairement un sommet :

\[
K\neq v.
\]

Elle peut être une sous-structure :

\[
K\equiv G_K\subseteq G_\theta.
\]

Cette représentation est cohérente avec l'hypothèse selon laquelle la
connaissance est une organisation relationnelle plutôt qu'un objet localisé.

% ============================================================
\section{Intelligence émergente}
% ============================================================

On définit un opérateur général de transformation :

\[
\Gamma_\theta:
\mathcal I
\rightarrow
\mathcal I,
\]

où $\mathcal I$ désigne l'espace des états informationnels internes.

L'intelligence fonctionnelle du système est alors associée non seulement à
la quantité d'information mémorisée, mais à la structure des transformations
qu'il peut effectuer :

\[
I
\rightarrow
\Gamma_\theta(I)
\rightarrow
\Gamma_\theta^2(I)
\rightarrow
\cdots.
\]

Ainsi :

\[
\boxed{
\text{intelligence}
\not\equiv
\text{quantité brute de mémoire}.
}
\]

Elle dépend également de la capacité du système à transformer et exploiter
les structures mémorisées.

% ============================================================
\section{Observation, mémoire, connaissance et intelligence}
% ============================================================

Le modèle complet peut être résumé par la chaîne :

\[
\boxed{
\begin{array}{ccccccccc}
\text{Domaine}
&
\xrightarrow{\text{observation}}
&
\text{Empreintes}
&
\xrightarrow{\text{persistance}}
&
\text{Mémoire}
&
\xrightarrow{\text{organisation}}
&
\text{Connaissance}
&
\xrightarrow{\text{transformation}}
&
\text{Intelligence}.
\end{array}
}
\]

Dans un LLM :

\[
\boxed{
\begin{array}{ccccccccc}
\D
&
\longrightarrow
&
I_\theta(\D)
&
\longrightarrow
&
\theta^\ast
&
\longrightarrow
&
\mathfrak M_{\theta^\ast}
&
\longrightarrow
&
\Gamma_{\theta^\ast}.
\end{array}
}
\]

Cette correspondance est une hypothèse de modélisation et non une identité
mathématique automatique entre mémoire humaine et paramètres d'un LLM.

% ============================================================
\section{Hypothèse centrale sur la boîte noire}
% ============================================================

\begin{hypothesis}[Hypothèse de la boîte noire par maillage d'empreintes]
La difficulté à interpréter un LLM pourrait provenir en partie du fait que
les significations fonctionnelles pertinentes ne sont pas localisées dans
des unités individuelles mais distribuées dans un maillage de structures
d'empreintes qui se recouvrent, se transforment et interagissent à plusieurs
niveaux de résolution.
\end{hypothesis}

Cette hypothèse peut être représentée par :

\[
\boxed{
\text{boîte noire}
\approx
\text{structure distribuée difficile à factoriser localement}.
}
\]

Elle conduit à remplacer la question :

\[
\text{``Que représente ce neurone ?''}
\]

par :

\[
\boxed{
\text{``Dans quelles empreintes et quelles structures relationnelles
ce neurone intervient-il ?''}
}
\]

% ============================================================
\section{Hypothèse de superposition}
% ============================================================

\begin{hypothesis}[Superposition structurelle]
Lorsque plusieurs structures informationnelles utilisent une même
sous-structure interne, leurs supports peuvent se recouvrir sans que les
structures globales soient identiques.
\end{hypothesis}

Formellement :

\[
R_a\neq R_b,
\qquad
R_a\cap R_b\neq\varnothing.
\]

Cette situation constitue le modèle théorique minimal de la poly-sémanticité.

Une unité $n$ peut satisfaire :

\[
n\in R_a
\]

et :

\[
n\in R_b
\]

tout en ayant :

\[
R_a\neq R_b.
\]

% ============================================================
\section{Hypothèse de généralisation par invariants}
% ============================================================

Soit une famille :

\[
\{R_i\}_{i\in I}.
\]

On recherche une structure :

\[
S
\]

telle que :

\[
S\subseteq R_i
\]

pour un nombre suffisamment important de $i$.

La généralisation est alors associée à la conservation de $S$ lorsque le
système rencontre une nouvelle donnée :

\[
R_{\mathrm{new}}
\]

telle que :

\[
S\subseteq R_{\mathrm{new}}.
\]

On obtient :

\[
\boxed{
\text{généralisation}
=
\text{réutilisation d'une structure relationnelle commune}.
}
\]

Cette définition permet de distinguer :

\[
\text{généralisation}
\]

de :

\[
\text{mémorisation littérale}.
\]

% ============================================================
\section{Hypothèse de résolution}
% ============================================================

La théorie propose que la représentation d'un même domaine puisse être
observée à plusieurs niveaux :

\[
\lambda_1<\lambda_2<\cdots<\lambda_m.
\]

À faible résolution, plusieurs états peuvent appartenir à la même classe :

\[
x\sim_{\lambda_1}y.
\]

À résolution supérieure :

\[
x\not\sim_{\lambda_2}y.
\]

Une représentation générale peut donc coexister avec des représentations
plus spécialisées.

Schématiquement :

\[
\boxed{
\text{général}
\leftarrow
\text{basse résolution}
}
\]

et :

\[
\boxed{
\text{spécialisation}
\leftarrow
\text{augmentation de résolution}.
}
\]

Cette formulation fournit une interprétation possible du passage de
catégories générales à des concepts plus précis.

% ============================================================
\section{Architecture conceptuelle complète}
% ============================================================

Le modèle complet peut être représenté comme suit :

\[
\begin{array}{cccccc}
\text{Dataset}
&
\longrightarrow
&
\text{Observations}
&
\longrightarrow
&
\text{Empreintes}
&
\\[0.2cm]
&&
\downarrow
&&
\downarrow
&
\\[0.2cm]
&&
\text{Dimensions}
&
\longrightarrow
&
\text{Interactions}
&
\\[0.2cm]
&&
&&
\downarrow
&
\\[0.2cm]
&&
&&
\text{Structures relationnelles}
&
\\[0.2cm]
&&
&&
\downarrow
&
\\[0.2cm]
&&
&&
\text{Maillage}
&
\\[0.2cm]
&&
&&
\downarrow
&
\\[0.2cm]
&&
&&
\text{Superposition}
&
\\[0.2cm]
&&
&&
\downarrow
&
\\[0.2cm]
&&
&&
\text{Généralisation}
&
\\[0.2cm]
&&
&&
\downarrow
&
\\[0.2cm]
&&
&&
\text{Comportement intelligent}.
&
\end{array}
\]

% ============================================================
\section{Correspondance avec la théorie générale des empreintes}
% ============================================================

La correspondance proposée est :

\begin{center}
\begin{tabular}{>{\raggedright}p{0.32\textwidth}
                >{\raggedright\arraybackslash}p{0.55\textwidth}}
\toprule
\textbf{Théorie des empreintes}
&
\textbf{Interprétation LLM proposée}
\\
\midrule
Observation
&
Traitement d'une donnée par le modèle
\\
Empreinte
&
État ou structure interne produite par le traitement
\\
Résolution
&
Niveau de détail auquel une représentation est observée
\\
Indiscernabilité
&
Équivalence de représentations relativement à une observation interne
\\
Classe d'équivalence
&
Ensemble de données indistinguables à une résolution donnée
\\
Seuil
&
Niveau auquel deux structures deviennent distinguables
\\
Quotient
&
Espace des représentations distinguées par une famille d'observations
\\
Fermeture
&
Ensemble d'informations déductibles à partir d'une empreinte donnée
\\
Transport
&
Transformation d'une représentation entre espaces ou niveaux
\\
Structure hiérarchique
&
Organisation des représentations à plusieurs niveaux de résolution
\\
\bottomrule
\end{tabular}
\end{center}

% ============================================================
\section{Précaution sur le mot ``émergence''}
% ============================================================

Le terme ``émergence'' doit être employé avec précision.

Le fait qu'une propriété soit difficile à attribuer à une unité individuelle
ne démontre pas automatiquement qu'elle est émergente au sens fort.

Dans le présent cadre, on utilisera donc :

\begin{definition}[Émergence structurelle]
Une propriété est dite structurellement émergente lorsqu'elle est définie
sur une configuration collective d'éléments et ne peut pas être identifiée
à la propriété isolée d'un élément individuel sans perdre une partie de sa
définition.
\end{definition}

Ainsi, si :

\[
P=P(R)
\]

où :

\[
R=\{n_1,\ldots,n_k\},
\]

et qu'il n'existe pas de fonction individuelle $p$ telle que :

\[
P(R)=p(n_i)
\]

pour un seul $i$, alors $P$ est une propriété structurelle du réseau.

% ============================================================
\section{Prédictions expérimentales}
% ============================================================

La théorie devient scientifiquement testable lorsqu'elle produit des
prédictions discriminantes.

\subsection{Prédiction 1 : non-localité de la signification}

Si une propriété sémantique $S$ correspond à une empreinte distribuée, alors
la suppression d'une seule unité ne devrait pas nécessairement supprimer
$S$.

On peut tester :

\[
\Delta S
=
S(G_\theta)-S(G_\theta\setminus n).
\]

Une faible valeur de $\Delta S$ pour certaines unités serait compatible
avec une représentation distribuée.

\subsection{Prédiction 2 : structure de recouvrement}

Pour deux familles sémantiques $A$ et $B$, on peut mesurer :

\[
J(A,B)
=
\frac{
|\mathcal C(A)\cap\mathcal C(B)|
}{
|\mathcal C(A)\cup\mathcal C(B)|
}.
\]

La théorie prédit qu'un recouvrement non nul ne suffit pas à établir
une identité sémantique.

\subsection{Prédiction 3 : dépendance au contexte}

Pour un même token $t$ :

\[
I_\theta(t\mid d_1)
\neq
I_\theta(t\mid d_2)
\]

dans de nombreux contextes.

Une analyse pertinente doit donc mesurer la variation :

\[
\Delta_I
=
d
\left(
I_\theta(t\mid d_1),
I_\theta(t\mid d_2)
\right).
\]

\subsection{Prédiction 4 : raffinement avec la résolution}

Si le modèle possède effectivement une organisation hiérarchique de
représentations, on devrait pouvoir observer des relations du type :

\[
\Pi_{\lambda_2}
\preceq
\Pi_{\lambda_1}
\]

où $\Pi_\lambda$ représente la partition induite par la résolution
$\lambda$.

\subsection{Prédiction 5 : généralisation par structures communes}

Une structure identifiée dans plusieurs familles de données devrait pouvoir
être réactivée dans une nouvelle famille présentant des relations
structurelles analogues, même lorsque les éléments lexicaux sont différents.

% ============================================================
\section{Limites actuelles}
% ============================================================

Le présent formalisme ne démontre pas les affirmations suivantes :

\begin{enumerate}
    \item que tout LLM possède nécessairement la structure décrite ;
    \item que toute poly-sémanticité observée provient de ce mécanisme ;
    \item que la théorie des empreintes constitue l'explication complète de
    la mémoire humaine ;
    \item que l'intelligence humaine et celle d'un LLM sont identiques ;
    \item que l'Univers est littéralement une mémoire ;
    \item que toute généralisation est nécessairement une conséquence d'un
    seuil de résolution ;
    \item que la superposition structurelle est la seule cause de la
    poly-sémanticité.
\end{enumerate}

Ces propositions constituent des hypothèses ou des programmes de recherche.

% ============================================================
\section{Thèse centrale}
% ============================================================

La thèse centrale du modèle peut être condensée dans l'énoncé suivant :

\begin{quote}
\textbf{
Un système intelligent peut être compris comme un système qui acquiert,
conserve, organise et transforme des empreintes produites par ses
observations. Lorsque ces empreintes sont distribuées dans une architecture
à plusieurs niveaux et que plusieurs structures peuvent partager des
composantes communes, une même unité interne peut participer à plusieurs
structures sémantiques distinctes. La généralisation résulte alors de
l'exploitation de structures communes entre plusieurs empreintes, tandis que
la spécialisation correspond à l'apparition de distinctions supplémentaires
lorsque la résolution augmente.
}
\end{quote}

Sous forme symbolique :

\[
\boxed{
\text{Observation}
\rightarrow
\text{Empreinte}
\rightarrow
\text{Organisation}
\rightarrow
\text{Maillage}
\rightarrow
\text{Superposition}
\rightarrow
\text{Généralisation}
\rightarrow
\text{Transformation}
\rightarrow
\text{Intelligence}.
}
\]

La boîte noire d'un LLM pourrait alors être reformulée non comme l'existence
d'un contenu mystérieusement caché dans chaque neurone, mais comme le
problème inverse consistant à reconstruire, à partir des observations
accessibles, une structure distribuée d'empreintes dont les composantes sont
partagées, contextualisées, hiérarchisées et dépendantes de la résolution.

\[
\boxed{
\text{Boîte noire}
=
\text{problème de reconstruction d'une structure interne d'empreintes}.
}
\]

% ============================================================
\section{Programme mathématique ultérieur}
% ============================================================

La formalisation ouvre plusieurs problèmes qui doivent être traités
séparément.

\subsection{Problème A : définition exacte de l'empreinte neuronale}

Construire une application :

\[
I_\theta:
\D
\rightarrow
\mathcal L_\theta
\]

dont la nature soit définie indépendamment de l'interprétation intuitive
des activations.

\subsection{Problème B : définition exacte du maillage}

Construire :

\[
\mathfrak M_\theta(\D)
\]

comme un objet mathématique doté d'une structure précise de nœuds, d'arêtes,
de relations et éventuellement de poids.

\subsection{Problème C : définition de la superposition}

Définir une mesure :

\[
\operatorname{Overlap}(R_i,R_j)
\]

permettant de distinguer :

\[
\text{simple coactivation},
\]

de :

\[
\text{partage structurel}.
\]

\subsection{Problème D : définition de la résolution interne}

Construire une filtration :

\[
\mathcal I_{\lambda_1}
\subseteq
\mathcal I_{\lambda_2}
\subseteq
\cdots
\]

compatible avec les empreintes réellement mesurables dans un LLM.

\subsection{Problème E : lien entre structure et comportement}

Déterminer les conditions sous lesquelles une propriété structurelle :

\[
S\subseteq\mathfrak M_\theta
\]

produit effectivement une propriété comportementale observable :

\[
B(S).
\]

\subsection{Problème F : distinction entre mémorisation et généralisation}

Définir des critères permettant de distinguer :

\[
\text{récupération d'une structure mémorisée}
\]

de :

\[
\text{construction d'une structure généralisée}.
\]

% ============================================================
\section{Conclusion}
% ============================================================

Le cadre développé ici propose une interprétation unifiée dans laquelle
l'observation constitue le mécanisme fondamental d'acquisition
d'information, l'empreinte constitue la représentation produite par cette
observation, la mémoire correspond à la persistance de ces structures, la
connaissance à leur organisation relationnelle et l'intelligence à leur
exploitation dynamique.

Appliqué aux LLM, le modèle conduit à abandonner une identification trop
simple :

\[
\text{neurone}=\text{concept}.
\]

Il propose plutôt :

\[
\boxed{
\text{neurone}
\in
\text{structure d'empreinte}
}
\]

et :

\[
\boxed{
\text{concept ou connaissance}
\approx
\text{configuration relationnelle distribuée}.
}
\]

Une même unité peut alors appartenir à plusieurs structures :

\[
n\in R_a\cap R_b,
\qquad
R_a\neq R_b,
\]

sans contradiction.

Le phénomène apparemment paradoxal du neurone associé à des concepts
éloignés peut ainsi être reformulé comme un problème de superposition
structurelle et de résolution observationnelle.

La proposition fondamentale reste cependant une hypothèse scientifique :

\[
\boxed{
\text{La structure distribuée des empreintes pourrait constituer
une partie essentielle de l'explication de la boîte noire des LLM.}
}
\]

Sa valeur scientifique dépendra de la possibilité de construire les objets
$\mathfrak M_\theta$, $I_{\theta,\lambda}$ et
$\operatorname{Overlap}$ de manière opérationnelle, puis de vérifier
expérimentalement les prédictions qui en découlent.
    
\end{document}